% GENERATED FILE. Edit canonical lessons, then run npm run generate:books.
% canonical-source-hash: fnv1a64:319a52dae8dfccb7
% canonical-lessons: ZH-C177-tuijian, ZH-C177-doufu, ZH-C177-bao, ZH-C177-xinxian, ZH-C177-nanchi

\chapter{Full Already}
\label{ch:zh-full-already}
\hlchaptermodality{177}

\begin{chapteropening}
\textbf{By the end of this chapter:} \emph{I can ask what to order, ask for tofu, say I am full, and say whether food is fresh or tastes bad.}

Complete the last lesson of chapter 177: I can ask what to order, ask for tofu, say I am full, and say whether food is fresh or tastes bad.
\end{chapteropening}

\section[tuījiàn]{tuījiàn — to recommend}

\label{lesson:ZH-C177-tuijian}



\begin{quote}
Before the new one: say the Mandarin for to topple over, then the Mandarin for to catch fire.
\end{quote}

\subsection*{You'll want to know: tuījiàn}
\emph{tuījiàn} — \textquotedblleft{}to recommend\textquotedblright{}.

\begin{quote}
\emph{fúwùyuán nǐ tuījiàn shénme cài} — Waiter, which dish do you recommend?
\end{quote}

Its characters come later; for now it is a word for the ear.

\begin{cousinweb}
\emph{tuījiàn} is built from \emph{tuī} 'to push' and \emph{jiàn} 'to put forward', both inherited from Middle Chinese, \emph{thwoj} and \emph{tsenH} (Guangyun, in Baxter's transcription), which glosses the second 'to present' and 'to put forward'.
\end{cousinweb}

\begin{grammarlens}[title={What you've built}]
One more word for asking for what I want.
\end{grammarlens}

\subsection*{Your turn}
\begin{itemize}
\raggedright
  \item \emph{Say it:} \emph{tuījiàn}
  \item \emph{Say it:} \emph{tuījiàn}, once more
  \item \emph{Say it:} \emph{fúwùyuán nǐ tuījiàn shénme cài} — Waiter, which dish do you recommend?
  \item \emph{Recall:} say the Mandarin for to topple over, then the Mandarin for to catch fire, then say \emph{tuījiàn} again
\end{itemize}

\begin{tcolorbox}[breakable,colback=teal!4,colframe=teal!35!black,title={Before you move on}]
What does tuījiàn mean? (\textquotedblleft{}To recommend\textquotedblright{}.) Say it once more.
\end{tcolorbox}
\section[dòufu]{dòufu — tofu; bean curd}

\label{lesson:ZH-C177-doufu}



\begin{quote}
Before the new one: say the Mandarin for to catch fire, then the Mandarin for to recommend.
\end{quote}

\subsection*{You'll want to know: dòufu}
\emph{dòufu} — \textquotedblleft{}tofu; bean curd\textquotedblright{}.

\begin{quote}
\emph{wǒmen yào yí fèn dòufu} — One portion of tofu for us, please.
\end{quote}

Its characters come later; for now it is a word for the ear.

\begin{cousinweb}
\emph{dòufu} is built from \emph{dòu} 'beans' and \emph{fu}, both inherited from Middle Chinese, \emph{duwH} and \emph{bjuX} (Guangyun, in Baxter's transcription), which glosses the second 'rotten' and 'broken down': beans broken down into curd. In the word the second loses its tone.
\end{cousinweb}

\begin{grammarlens}[title={What you've built}]
One more word for asking for what I want.
\end{grammarlens}

\subsection*{Your turn}
\begin{itemize}
\raggedright
  \item \emph{Say it:} \emph{dòufu}
  \item \emph{Say it:} \emph{dòufu}, once more
  \item \emph{Say it:} \emph{wǒmen yào yí fèn dòufu} — One portion of tofu for us, please.
  \item \emph{Recall:} say the Mandarin for to catch fire, then the Mandarin for to recommend, then say \emph{dòufu} again
\end{itemize}

\begin{tcolorbox}[breakable,colback=teal!4,colframe=teal!35!black,title={Before you move on}]
What does dòufu mean? (\textquotedblleft{}Tofu; bean curd\textquotedblright{}.) Say it once more.
\end{tcolorbox}
\section[bǎo]{bǎo — full after eating; full up}

\label{lesson:ZH-C177-bao}



\begin{quote}
Before the new one: say the Mandarin for to recommend, then the Mandarin for tofu.
\end{quote}

\subsection*{You'll want to know: bǎo}
\emph{bǎo} — \textquotedblleft{}full after eating; full up\textquotedblright{}.

\begin{quote}
\emph{wǒ hěn bǎo bù chī le} — I'm full: no more for me.
\end{quote}

Its characters come later; for now it is a word for the ear.

\begin{cousinweb}
\emph{bǎo} is inherited from Middle Chinese \emph{paewX} (Guangyun, in Baxter's transcription), glossed 'to have eaten much'.
\end{cousinweb}

\begin{grammarlens}[title={What you've built}]
One more word for asking for what I want.
\end{grammarlens}

\subsection*{Your turn}
\begin{itemize}
\raggedright
  \item \emph{Say it:} \emph{bǎo}
  \item \emph{Say it:} \emph{bǎo}, once more
  \item \emph{Say it:} \emph{wǒ hěn bǎo bù chī le} — I'm full: no more for me.
  \item \emph{Recall:} say the Mandarin for to recommend, then the Mandarin for tofu, then say \emph{bǎo} again
\end{itemize}

\begin{tcolorbox}[breakable,colback=teal!4,colframe=teal!35!black,title={Before you move on}]
What does bǎo mean? (\textquotedblleft{}Full after eating; full up\textquotedblright{}.) Say it once more.
\end{tcolorbox}
\section[xīnxiān]{xīnxiān — fresh}

\label{lesson:ZH-C177-xinxian}



\begin{quote}
Before the new one: say the Mandarin for tofu, then the Mandarin for full after eating.
\end{quote}

\subsection*{You'll want to know: xīnxiān}
\emph{xīnxiān} — \textquotedblleft{}fresh\textquotedblright{}.

\begin{quote}
\emph{zhèr de yú gèng xīnxiān} — The fish here is fresher.
\end{quote}

Its characters come later; for now it is a word for the ear.

\begin{cousinweb}
\emph{xīnxiān} is built from \emph{xīn} 'new' and \emph{xiān} 'bright and clean', both inherited from Middle Chinese, \emph{sin} and \emph{sjen} (Guangyun, in Baxter's transcription).
\end{cousinweb}

\begin{grammarlens}[title={What you've built}]
One more word for asking for what I want.
\end{grammarlens}

\subsection*{Your turn}
\begin{itemize}
\raggedright
  \item \emph{Say it:} \emph{xīnxiān}
  \item \emph{Say it:} \emph{xīnxiān}, once more
  \item \emph{Say it:} \emph{zhèr de yú gèng xīnxiān} — The fish here is fresher.
  \item \emph{Recall:} say the Mandarin for tofu, then the Mandarin for full after eating, then say \emph{xīnxiān} again
\end{itemize}

\begin{tcolorbox}[breakable,colback=teal!4,colframe=teal!35!black,title={Before you move on}]
What does xīnxiān mean? (\textquotedblleft{}Fresh\textquotedblright{}.) Say it once more.
\end{tcolorbox}
\section[nánchī]{nánchī — tasting bad; horrible to eat}

\label{lesson:ZH-C177-nanchi}



\begin{quote}
Before the new one: say the Mandarin for full after eating, then the Mandarin for fresh.
\end{quote}

\subsection*{You'll want to know: nánchī}
\emph{nánchī} — \textquotedblleft{}tasting bad; horrible to eat\textquotedblright{}.

\begin{quote}
\emph{zhège cài bù xīnxiān hěn nánchī} — This dish isn't fresh, and it tastes horrible.
\end{quote}

Its characters come later; for now it is a word for the ear.

\begin{cousinweb}
\emph{nánchī} is built from \emph{nán} 'hard' and \zh{吃} 'to eat': hard to eat. The first is inherited from Middle Chinese \emph{nan} (Guangyun, in Baxter's transcription); for the second the dictionary has \emph{khek}, glossed 'to eat', under another way of writing it, though \emph{khek} would regularly give qì.
\end{cousinweb}

\begin{grammarlens}[title={What you've built}]
One more word for asking for what I want.
\end{grammarlens}

\subsection*{Your turn}
\begin{itemize}
\raggedright
  \item \emph{Say it:} \emph{nánchī}
  \item \emph{Say it:} \emph{nánchī}, once more
  \item \emph{Say it:} \emph{zhège cài bù xīnxiān hěn nánchī} — This dish isn't fresh, and it tastes horrible.
  \item \emph{Recall:} say the Mandarin for full after eating, then the Mandarin for fresh, then say \emph{nánchī} again
\end{itemize}

\begin{tcolorbox}[breakable,colback=teal!4,colframe=teal!35!black,title={Before you move on}]
What does nánchī mean? (\textquotedblleft{}Tasting bad; horrible to eat\textquotedblright{}.) Say it once more.
\end{tcolorbox}
